% !TeX program = xelatex
\documentclass[a4paper,10pt]{article}

% ============================================================================
% សន្លឹកកិច្ចការសិស្ស ស-ខ្មៅ (Student Worksheet - Grade 8 3.2.1 & 3.2.2 PCT1)
% បរិបទ៖ គ្មានស្លាយ LCD និងគ្មានម៉ាស៊ីនគិតលេខ — សិស្សបំពេញផ្ទាល់លើសន្លឹកកិច្ចការ ២ ទំព័រ A4
% ============================================================================
\usepackage[a4paper, left=0.85cm, right=0.85cm, top=0.70cm, bottom=0.95cm, footskip=0.48cm]{geometry}
\usepackage{amsmath, amssymb}
\usepackage{unicode-math}
\usepackage{fontspec}
\usepackage{array, colortbl}
\usepackage[table]{xcolor}
\usepackage{enumitem}
\usepackage{pifont}
\usepackage{fancyhdr}
\usepackage{lastpage}
\usepackage{tikz}
\usepackage[most]{tcolorbox}
\usetikzlibrary{calc, arrows.meta, decorations.pathreplacing, positioning, babel}

\makeatletter
\let\@expl@@@mark@update@singlecol@structures@@\relax
\makeatother

\providecommand{\NGSFindFonts}{%
  \IfFileExists{C:/Windows/Fonts/cambria.ttc}{%
    \def\NGSSysPath{C:/Windows/Fonts/}%
    \def\NGSUserPath{C:/Users/KIT/AppData/Local/Microsoft/Windows/Fonts/}%
  }{%
    \IfFileExists{./fonts/KhmerOSsiemreap.ttf}{\def\NGSSysPath{./fonts/}\def\NGSUserPath{./fonts/}}{%
    \IfFileExists{../fonts/KhmerOSsiemreap.ttf}{\def\NGSSysPath{../fonts/}\def\NGSUserPath{../fonts/}}{%
    \IfFileExists{../../fonts/KhmerOSsiemreap.ttf}{\def\NGSSysPath{../../fonts/}\def\NGSUserPath{../../fonts/}}{%
    \IfFileExists{../../../fonts/KhmerOSsiemreap.ttf}{\def\NGSSysPath{../../../fonts/}\def\NGSUserPath{../../../fonts/}}{%
    \IfFileExists{../../../../fonts/KhmerOSsiemreap.ttf}{\def\NGSSysPath{../../../../fonts/}\def\NGSUserPath{../../../../fonts/}}{%
    \IfFileExists{../../../../../fonts/KhmerOSsiemreap.ttf}{\def\NGSSysPath{../../../../../fonts/}\def\NGSUserPath{../../../../../fonts/}}{%
    \IfFileExists{../../../../../../fonts/KhmerOSsiemreap.ttf}{\def\NGSSysPath{../../../../../../fonts/}\def\NGSUserPath{../../../../../../fonts/}}{%
      \def\NGSSysPath{../../../../../../../fonts/}\def\NGSUserPath{../../../../../../../fonts/}%
    }}}}}}}%
  }%
}
\NGSFindFonts

\setmathfont[Path={\NGSSysPath}, FontIndex=1, Scale=1.02]{cambria.ttc}
\defaultfontfeatures{Ligatures=TeX}

\setmainfont[
  Path={\NGSUserPath}, Extension=.ttf, Script=Khmer,
  BoldFont={KhmerOSsiemreap}, BoldFeatures={FakeBold=1.8},
  ItalicFont={KhmerOSsiemreap}, ItalicFeatures={FakeSlant=0.22}
]{KhmerOSsiemreap}

\newfontfamily\muol[Path={\NGSUserPath}, Extension=.ttf, Script=Khmer]{KhmerOSmuollight}
\newfontfamily\TimesEn[Path={\NGSSysPath}, Extension=.ttf, BoldFont={timesbd}, ItalicFont={timesi}]{times}

\XeTeXlinebreaklocale "km"
\XeTeXlinebreakskip = 0pt plus 1pt minus 0.1pt

\definecolor{panelgray}{gray}{0.93}
\definecolor{headgray}{gray}{0.85}

\pagestyle{fancy}
\fancyhf{}
\renewcommand{\headrulewidth}{0pt}
\fancyfoot[L]{\fontsize{7.5pt}{9.5pt}\selectfont អនុវិទ្យាល័យព្រះក្រសាំង \en{|} គណិតវិទ្យាថ្នាក់ទី៨ \en{|} មេរៀនទី៣៖ ២.១ ភាគរយនៃបរិមាណ និង ២.២ ចំណេញ-ខាត (គ្រូ៖ ហុក ប៊ុនថែន)}
\fancyfoot[R]{\fontsize{7.8pt}{9.5pt}\selectfont ទំព័រ \en{\thepage/\pageref{LastPage}}}

\setlength{\parindent}{0pt}
\setlength{\parskip}{0pt}
\setlength{\tabcolsep}{3.8pt}

\newcommand{\en}[1]{{\TimesEn #1}}
\newcommand{\enbf}[1]{{\TimesEn\bfseries #1}}
\newcommand{\unit}[1]{\,\text{\en{#1}}}
\newcommand{\cnum}[1]{\mbox{\fontsize{10pt}{12pt}\selectfont\ding{\numexpr181+#1\relax}}}
\newcommand{\chkbox}{\tikz[baseline=-0.12ex]\draw[black, line width=0.8pt] (0,0) rectangle (0.24cm,0.24cm);}
\newcommand{\blank}[1][1.5cm]{\makebox[#1]{\dotfill}}
\newcommand{\ansbox}[1][0.85cm]{\,\tikz[baseline=-0.5ex]\draw[black, line width=0.65pt, fill=white] (0,0) rectangle (#1, 0.42cm);\,}
\newcommand{\secbar}[2]{%
  \begin{tcolorbox}[enhanced, sharp corners, colback=headgray, colframe=black, boxrule=0.9pt, boxsep=1pt, top=2pt, bottom=2pt, left=5pt, right=5pt, before skip=4pt, after skip=3pt]
    {\muol\fontsize{9.2pt}{13pt}\selectfont #1}\hfill{\fontsize{8pt}{11pt}\selectfont\textbf{#2}}
  \end{tcolorbox}%
}

\newcolumntype{L}[1]{>{\raggedright\arraybackslash}p{#1}}
\newcolumntype{C}[1]{>{\centering\arraybackslash}p{#1}}

\begin{document}
\fontsize{9.6pt}{14.4pt}\selectfont

% ============================================================================
% ទំព័រទី ១៖ E1 + E2 (POE 1 & POE 2) + E3 សំយោគរូបមន្ត និងធនាគារហេតុផល
% ============================================================================
\begin{tcolorbox}[enhanced, sharp corners, colback=white, colframe=black, boxrule=1.0pt, boxsep=2pt, top=3pt, bottom=3pt, left=6pt, right=6pt]
  \begin{center}
    {\muol\fontsize{10.1pt}{15pt}\selectfont អនុវិទ្យាល័យព្រះក្រសាំង \enbf{--} សន្លឹកកិច្ចការសិស្ស \enbf{(}មេរៀនទី៣៖ ២\enbf{.}១\enbf{--}២\enbf{.}២ ភាគរយ និងចំណេញ-ខាត\enbf{)}}\\[3pt]
    {\fontsize{8.8pt}{12pt}\selectfont
    ឈ្មោះសិស្ស / ដៃគូ៖ \blank[5.2cm] \hspace{0.25cm}
    ថ្នាក់ទី៨៖ \blank[1.1cm] \hspace{0.25cm}
    ក្រុមទី៖ \blank[1.0cm] \hspace{0.25cm}
    កាលបរិច្ឆេទ៖ ២៨/១១/២០២៦}
  \end{center}
\end{tcolorbox}

% --- E1 + E2 ---
\secbar{ផ្នែកទី១៖ បេសកកម្មទស្សន៍ទាយ \enbf{(E1)} និងការរុករកជាដៃគូ \enbf{(E2 -- POE 1 \& POE 2)}}{រយៈពេល ១៤ នាទី}

\begin{tcolorbox}[enhanced, sharp corners, colback=panelgray, colframe=black, boxrule=0.75pt, boxsep=1.5pt, top=2.5pt, bottom=2.5pt, left=5pt, right=5pt]
  \textbf{បេសកកម្ម \en{E1} (គូសធីក \chkbox\ មុនពិសោធន៍)៖}
  (១) ចម្ងាយ $45\unit{m}$ ធៀបនឹង $1\unit{km}$ ស្មើប៉ុន្មានភាគរយ? \chkbox\ $45\%$ \quad \chkbox\ $4.5\%$\\
  (២) ទិញអាវថ្លៃ $20\,000\text{៛}$ រួចលក់ចេញថ្លៃ $25\,000\text{៛}$ — តើចំណេញប៉ុន្មានភាគរយ? \chkbox\ $20\%$ \quad \chkbox\ $25\%$
\end{tcolorbox}

\vspace{3pt}
\textbf{ក. សកម្មភាពទី ១ (\en{POE 1})៖ ច្រកឯកតារួមលើភាគរយនៃបរិមាណមួយ (សៀវភៅពុម្ព ទំ.៣១)}

{\renewcommand{\arraystretch}{1.46}
\begin{tabular}{|L{4.30cm}|L{4.80cm}|L{5.30cm}|C{3.20cm}|}
  \hline
  \rowcolor{headgray}
  \textbf{បរិមាណទាំងពីរ} & \textbf{ជំហានប្ដូរជាឯកតារួម} & \textbf{ប្រភាគគុណនឹង $100\%$} & \textbf{លទ្ធផល ($\%$)} \\
  \hline
  $45\unit{m}$ នៃ $1\unit{km}$ & $1\unit{km} = \blank[1.3cm]\unit{m}$ & $\dfrac{45}{\blank[1.2cm]} \times 100\%$ & $\blank[1.4cm]\%$ \\
  \hline
  សិស្សស្រី $30$ នៃសិស្ស $40$ នាក់ & ឯកតាដូចគ្នា (នាក់) & $\dfrac{30}{40} \times 100\% = \dfrac{3}{4}\times 100\%$ & $\blank[1.4cm]\%$ \\
  \hline
  $250\unit{g}$ នៃ $2\unit{kg}$ & $2\unit{kg} = \blank[1.3cm]\unit{g}$ & $\dfrac{250}{\blank[1.2cm]} \times 100\%$ & $\blank[1.4cm]\%$ \\
  \hline
  $15\unit{mn}$ នៃ $1\unit{h}$ & $1\unit{h} = \blank[1.3cm]\unit{mn}$ & $\dfrac{15}{\blank[1.2cm]} \times 100\%$ & $\blank[1.4cm]\%$ \\
  \hline
\end{tabular}}

\vspace{3pt}
$\bullet$ \textbf{ពន្យល់ (\en{E})៖} មុនសរសេរបរិមាណ $A$ ជាភាគរយនៃបរិមាណ $B$ យើងត្រូវប្ដូរបរិមាណទាំងពីរឱ្យមាន \blank[2.3cm] ជាមុនសិន!

\vspace{4pt}
\textbf{ខ. សកម្មភាពទី ២ (\en{POE 2})៖ របារយុថ្កាគោលថ្លៃដើម $100\%$ — ប្រៀបធៀបចំណេញ និងខាត (ទំ.៣២)}

{\renewcommand{\arraystretch}{1.46}
\begin{tabular}{|L{4.80cm}|C{2.70cm}|C{2.70cm}|C{3.40cm}|C{4.00cm}|}
  \hline
  \rowcolor{headgray}
  \textbf{ស្ថានភាពទិញ-លក់} & \textbf{ថ្លៃដើម ($100\%$)} & \textbf{ថ្លៃលក់} & \textbf{ចំណេញ ឬ ខាត (៛)} & \textbf{ភាគរយធៀបនឹងថ្លៃដើម} \\
  \hline
  ១) ទិញអាវ $20\,000\text{៛}$ លក់ $25\,000\text{៛}$ & $20\,000\text{៛}$ & $25\,000\text{៛}$ & ចំណេញ $\blank[1.2cm]$ & $\dfrac{5000}{\blank[1.2cm]}\times 100\% = \blank[0.9cm]\%$ \\
  \hline
  ២) ទិញអាវ $25\,000\text{៛}$ លក់ $20\,000\text{៛}$ & $25\,000\text{៛}$ & $20\,000\text{៛}$ & ខាត $\blank[1.2cm]$ & $\dfrac{5000}{\blank[1.2cm]}\times 100\% = \blank[0.9cm]\%$ \\
  \hline
\end{tabular}}

\vspace{3pt}
$\bullet$ \textbf{ពន្យល់ (\en{E})៖} ហេតុអ្វីប្រាក់ខុសគ្នា $5000\text{៛}$ ដូចគ្នា តែភាគរយមិនស្មើគ្នា ($25\% \ne 20\%$)? ព្រោះភាគបែងត្រូវយក \blank[2.4cm] ជានិច្ច!

% --- E3 ---
\secbar{ផ្នែកទី២៖ សំយោគរូបមន្តគោល \enbf{(E3)} និងធនាគារហេតុផល \enbf{B\&W} \cnum{1}\enbf{--}\cnum{5}}{សៀវភៅពុម្ព ទំ.៣១–៣២}

\begin{tabular}{|L{9.20cm}|L{9.20cm}|}
  \hline
  \textbf{១. រូបមន្តគោល (ទំ.៣១–៣២)៖}\par\vspace{2pt}
  $\bullet$ \textbf{ភាគរយនៃបរិមាណ (ឯកតាដូចគ្នា)៖} $\dfrac{A}{B} \times \blank[1.2cm]\%$\par\vspace{3pt}
  $\bullet$ \textbf{ភាគរយចំណេញ៖} $\dfrac{\text{ប្រាក់ចំណេញ}}{\blank[2.2cm]} \times 100\%$\par\vspace{3pt}
  $\bullet$ \textbf{ភាគរយខាត៖} $\dfrac{\text{ប្រាក់ខាត}}{\blank[2.2cm]} \times 100\%$\par\vspace{2pt}
  $\bullet$ \textbf{ថ្លៃលក់៖} ចំណេញ $p\% \Rightarrow \text{ដើម}\!\times\!\frac{100+p}{100}$ ; ខាត $q\% \Rightarrow \text{ដើម}\!\times\!\frac{100-q}{100}$
  &
  \textbf{២. ធនាគារហេតុផល \en{B\&W} (\cnum{1}--\cnum{5})៖}\par\vspace{2pt}
  \cnum{1} ច្រកឯកតារួម៖ ប្ដូរបរិមាណទាំងពីរឱ្យមានឯកតាដូចគ្នាមុនចែក ។\par\vspace{2pt}
  \cnum{2} យុថ្កាថ្លៃដើម៖ ថ្លៃដើមស្មើ $100\%$ ជាភាគបែងជានិច្ច (មិនមែនថ្លៃលក់ទេ!) ។\par\vspace{2pt}
  \cnum{3} ចំណេញ $p\%$៖ ថ្លៃលក់ស្មើ $(100+p)\%$ នៃប្រាក់ថ្លៃដើម ។\par\vspace{2pt}
  \cnum{4} ខាត $q\%$៖ ថ្លៃលក់ស្មើ $(100-q)\%$ នៃប្រាក់ថ្លៃដើម ។\par\vspace{2pt}
  \cnum{5} លក់បន្តពីរដំណាក់៖ ថ្លៃទិញរបស់ម្ចាស់ថ្មី គឺជាថ្លៃដើមថ្មី ($100\%$) ។ \\
  \hline
\end{tabular}

\vspace{3pt}
\begin{tcolorbox}[enhanced, sharp corners, colback=panelgray, colframe=black, boxrule=0.65pt, boxsep=1.5pt, top=2pt, bottom=2pt, left=5pt, right=5pt]
  \textbf{ប្រយ័ត្នអន្ទាក់៖} (១) មុនគណនាភាគរយ ត្រូវប្ដូរបរិមាណទាំងពីរឱ្យមាន\textbf{ឯកតាដូចគ្នា}! \quad (២) ភាគរយចំណេញ និងខាតត្រូវចែកនឹង\textbf{ថ្លៃដើម}ជានិច្ច!
\end{tcolorbox}

\vspace{3pt}
\textbf{៣. លំហាត់ត្រួតពិនិត្យរហ័ស៖ គណនាភាគរយនៃបរិមាណ និងភាគរយចំណេញ-ខាត}

\vspace{2pt}
{\renewcommand{\arraystretch}{1.60}
\begin{tabular}{|C{4.40cm}|C{4.40cm}|C{4.40cm}|C{4.40cm}|}
  \hline
  \rowcolor{panelgray}
  $20\%$ នៃ $150\,000\text{៛} = \blank[1.0cm]$ & $12$ នាក់នៃ $50$ នាក់ $= \blank[0.8cm]\%$ & $80\unit{cm}$ នៃ $4\unit{m} = \blank[0.8cm]\%$ & $300\unit{m}$ នៃ $2\unit{km} = \blank[0.8cm]\%$ \\
  \hline
  ទិញ $50\$$ លក់ $60\$$៖ ចំណេញ $\blank[0.7cm]\%$ & ទិញ $80\$$ លក់ $68\$$៖ ខាត $\blank[0.7cm]\%$ & ដើម $40\,000\text{៛}$ ចំណេញ $10\% \to \blank[0.8cm]$ & ដើម $50\,000\text{៛}$ ខាត $20\% \to \blank[0.8cm]$ \\
  \hline
  $30\unit{mn}$ នៃ $2\unit{h} = \blank[0.8cm]\%$ & $400\unit{g}$ នៃ $1\unit{kg} = \blank[0.8cm]\%$ & ទិញ $200\$$ លក់ $230\$$៖ $+\blank[0.7cm]\%$ & ទិញ $120\$$ លក់ $90\$$៖ $-\blank[0.7cm]\%$ \\
  \hline
  $15\%$ នៃ $80\,000\text{៛} = \blank[1.0cm]$ & $18$ នាក់នៃ $60$ នាក់ $= \blank[0.8cm]\%$ & ដើម $30\,000\text{៛}$ ចំណេញ $25\% \to \blank[0.8cm]$ & ដើម $60\,000\text{៛}$ ខាត $15\% \to \blank[0.8cm]$ \\
  \hline
\end{tabular}}

\newpage

% ============================================================================
% ទំព័រទី ២៖ E4 លំហាត់គំរូ + E5 ប្រតិបត្តិ + វិនិច្ឆ័យ A-D + Exit Check + កិច្ចការផ្ទះ
% ============================================================================
\secbar{ផ្នែកទី៣៖ ពង្រីកចំណេះដឹងលើការលក់បន្តពីរដំណាក់កាល \enbf{(E4 --} គំរូ ទំ\enbf{.}៣២\enbf{)}}{រយៈពេល ១០ នាទី}

\textbf{លំហាត់គំរូ (សៀវភៅពុម្ព ទំ.៣២)៖} បុរសម្នាក់ទិញរថយន្តមួយគ្រឿងថ្លៃ $\mathbf{75\,000\$}$ រួចលក់ទៅឱ្យមិត្តភក្ដិគាត់វិញថ្លៃ $\mathbf{84\,000\$}$ ។ មួយឆ្នាំក្រោយមក មិត្តភក្ដិគាត់លក់រថយន្តនោះបន្តបានត្រឹមតែ $\mathbf{60\,000\$}$ ។ គណនាភាគរយចំណេញ និងភាគរយខាតក្នុងដំណាក់កាលនីមួយៗ៖

{\renewcommand{\arraystretch}{1.48}
\begin{tabular}{|L{4.50cm}|L{9.10cm}|C{4.20cm}|}
  \hline
  \rowcolor{headgray}
  \textbf{ដំណាក់កាលទិញ-លក់} & \textbf{ការកំណត់យុថ្កាថ្លៃដើម និងការគណនាដោយដៃ} & \textbf{ចម្លើយ និងហេតុផល} \\
  \hline
  \textbf{១. ម្ចាស់ទី ១ លក់ឱ្យមិត្តភក្ដិ} & ថ្លៃដើមទី១ $= \blank[1.3cm]\$$ ; ចំណេញ $= 84\,000 - 75\,000 = \blank[1.2cm]\$$ & យុថ្កាទី១ $= 75\,000\$$ (\cnum{2}) \\
  \hline
  \textbf{២. គណនាភាគរយចំណេញ} & $\text{ចំណេញ}\% = \dfrac{\blank[1.2cm]}{75\,000}\times 100\% = \dfrac{900}{75}\% = \blank[1.1cm]\%$ & ចំណេញ $\blank[0.9cm]\%$ (\cnum{2}) \\
  \hline
  \textbf{៣. មិត្តភក្ដិលក់បន្ត ($60\,000\$$)} & ថ្លៃដើមថ្មី $= \blank[1.3cm]\$$ ; ប្រាក់ខាត $= 84\,000 - 60\,000 = \blank[1.2cm]\$$ & យុថ្កាថ្មី $= 84\,000\$$ (\cnum{5}) \\
  \hline
  \textbf{៤. គណនាភាគរយខាត} & $\text{ខាត}\% = \dfrac{24\,000}{\blank[1.3cm]}\times 100\% = \dfrac{2}{7}\times 100\% = 28\dfrac{\blank[0.5cm]}{7}\% \approx \blank[1.0cm]\%$ & ខាត $\approx \blank[1.0cm]\%$ (\cnum{5}) \\
  \hline
\end{tabular}}

\vspace{4pt}
$\bullet$ \textbf{សំណួរឆ្លុះបញ្ចាំង \en{(EIA Looking-Back)}៖} ហេតុអ្វីនៅដំណាក់ទី ២ យើងមិនអាចយក $24\,000\$$ ចែកនឹង $75\,000\$$ បានទេ? \blank[3.0cm]

% --- E5 ---
\secbar{ផ្នែកទី៤៖ វាយតម្លៃលទ្ធផលសិក្សា \enbf{(E5 --} ប្រតិបត្តិ ១ និង ២ \enbf{+} វិនិច្ឆ័យ \enbf{A--D +} សំបុត្រចេញ\enbf{)}}{រយៈពេល ១០ នាទី}

\textbf{ក. លំហាត់ប្រតិបត្តិ (សៀវភៅពុម្ព ទំ.៣២ លើ និង ទំ.៣៣ លើ — ធ្វើជាដៃគូ រួចប្ដូរសន្លឹកផ្ទៀង \en{Pause-Review})៖}

{\renewcommand{\arraystretch}{1.48}
\begin{tabular}{|L{8.90cm}|L{8.90cm}|}
  \hline
  \textbf{១) ប្រតិបត្តិទី ១ (ទំ.៣២ — ប្រៀបធៀបពិន្ទុប្រឡង)៖}\par\vspace{2pt}
  សិស្សម្នាក់ប្រឡងគណិតវិទ្យាបាន $\dfrac{76}{80}$ ពិន្ទុ និងភាសាខ្មែរបាន $\dfrac{46}{50}$ ពិន្ទុ ។\par\vspace{2pt}
  $\bullet$ គណិតវិទ្យា៖ $\dfrac{76}{80}\times 100\% = \blank[1.4cm]\%$\par\vspace{2pt}
  $\bullet$ ភាសាខ្មែរ៖ $\dfrac{46}{50}\times 100\% = \blank[1.4cm]\%$\par\vspace{2pt}
  $\bullet$ មុខវិជ្ជាដែលបានភាគរយខ្ពស់ជាង៖ \chkbox\ គណិតវិទ្យា \quad \chkbox\ ភាសាខ្មែរ
  &
  \textbf{២) ប្រតិបត្តិទី ២ (ទំ.៣៣ — ម៉ូតូរបស់សុខ)៖}\par\vspace{2pt}
  សុខទិញម៉ូតូមួយគ្រឿងថ្លៃ $\mathbf{2\,240\,000\text{៛}}$ ។ ជិះបាន ១ ឆ្នាំ គាត់លក់វិញ\textbf{ខាត $22\%$} ។ រកថ្លៃលក់ម៉ូតូនោះ៖\par\vspace{2pt}
  $\bullet$ ប្រាក់ខាត $= 2\,240\,000 \times \dfrac{22}{100} = \blank[1.6cm]\text{៛}$\par\vspace{2pt}
  $\bullet$ ថ្លៃលក់ $= 2\,240\,000 - \blank[1.5cm] = \blank[1.8cm]\text{៛}$ (\cnum{4})\par\vspace{2pt}
  (ឬគណនាផ្ទាល់៖ $2\,240\,000 \times \dfrac{78}{100} = \blank[1.8cm]\text{៛}$) \\
  \hline
\end{tabular}}

\vspace{5pt}
\textbf{ខ. គូសវង់ជុំវិញជម្រើសត្រឹមត្រូវតែ ១ គត់ ក្នុងចំណោម \en{A, B, C, D} រួចសរសេរលេខហេតុផល \cnum{1}--\cnum{5} នៃជម្រើសខុស៖}

{\renewcommand{\arraystretch}{1.36}
\begin{tabular}{|C{4.40cm}|C{4.40cm}|C{4.40cm}|C{4.40cm}|}
  \hline
  \cellcolor{panelgray}\textbf{ជម្រើស \en{A}} & \cellcolor{panelgray}\textbf{ជម្រើស \en{B}} & \cellcolor{panelgray}\textbf{ជម្រើស \en{C}} & \cellcolor{panelgray}\textbf{ជម្រើស \en{D}} \\
  \hline
  ម៉ូតូ $2\,240\,000\text{៛}$ ខាត $22\%$ លក់បាន $1\,747\,200\text{៛}$ & $45\unit{m}$ នៃ $1\unit{km}$ ស្មើនឹង $45\%$ & ទិញ $20\,000\text{៛}$ លក់ $25\,000\text{៛}$ ចំណេញ $20\%$ & ពិន្ទុខ្មែរ $\frac{46}{50}$ ខ្ពស់ជាងគណិត $\frac{76}{80}$ \\[3pt]
  \chkbox\ ត្រូវ \quad \chkbox\ ខុស (\blank[0.7cm]) & \chkbox\ ត្រូវ \quad \chkbox\ ខុស (\blank[0.7cm]) & \chkbox\ ត្រូវ \quad \chkbox\ ខុស (\blank[0.7cm]) & \chkbox\ ត្រូវ \quad \chkbox\ ខុស (\blank[0.7cm]) \\
  \hline
\end{tabular}}

\vspace{5pt}
\begin{tcolorbox}[enhanced, sharp corners, colback=white, colframe=black, boxrule=0.85pt, boxsep=2.5pt, top=3.5pt, bottom=3.5pt, left=5pt, right=5pt]
  {\muol\fontsize{8.8pt}{12.5pt}\selectfont សំបុត្រចេញរហ័ស \enbf{(Exit Check --} ៦០ វិនាទី\enbf{)}៖}\\[2pt]
  (១) សរសេរ $300\unit{g}$ ជាភាគរយនៃ $2\unit{kg}$៖ $\dfrac{300}{\blank[1.1cm]}\times 100\% = \blank[1.1cm]\%$ \quad\en{|}\quad
  (២) ទិញកាបូប $40\,000\text{៛}$ លក់ចំណេញ $15\%$ តើលក់ថ្លៃប៉ុន្មាន? \blank[2.0cm]៛
\end{tcolorbox}

\vspace{3pt}
\begin{tcolorbox}[enhanced, sharp corners, colback=panelgray, colframe=black, boxrule=0.75pt, boxsep=2pt, top=3pt, bottom=3pt, left=5pt, right=5pt]
  {\muol\fontsize{8.8pt}{12.5pt}\selectfont កិច្ចការផ្ទះ \enbf{(}សៀវភៅគោល ទំព័រ ៤០\enbf{--}៤១ លំហាត់ \enbf{\#3(}ក\enbf{,}ខ\enbf{,}គ\enbf{,}ឃ\enbf{)} និង \enbf{\#5(}ខ\enbf{,}គ\enbf{))}៖}\\[2pt]
  $\bullet$ \textbf{លំហាត់ \#3 (ទំ.៤០)៖} សរសេរបរិមាណទី១ ជាភាគរយនៃបរិមាណទី២ ដោយប្ដូរឱ្យមានឯកតារួមជាមុន ។\\
  $\bullet$ \textbf{លំហាត់ \#5ខ-គ (ទំ.៤១)៖} គណនាភាគរយនៃបរិមាណ និងប្រាក់ចំណេញ-ខាត រួចអានមុន \textbf{២.៣. ការបញ្ចុះតម្លៃ និង ២.៤. បម្រែបម្រួលភាគរយ (ទំ.៣៣–៣៥)} ។
\end{tcolorbox}

\end{document}
